\documentclass[10pt,a4paper]{amsart}
\usepackage[margin=19mm]{geometry}
\usepackage{amsmath,amssymb,mathtools}
\usepackage[scr=rsfs,cal=euler]{mathalfa}
\usepackage{xfrac}
\usepackage[unicode,hidelinks,bookmarks=false]{hyperref}
\newcommand{\ie}{i.e.\ }
\newcommand{\cf}{cf.\ }
\newcommand{\resp}{respectively\ }
\newcommand{\Subsection}{\subsection}
\providecommand{\nobreakdash}{\nobreak\hbox{-}}
\setlength{\emergencystretch}{2em}
\multlinegap0pt
\usepackage{bm}
\usepackage{bbm}
\usepackage{dsfont}
\usepackage{xspace}
\usepackage{cancel}
\usepackage{mathtools}
\usepackage{amsthm}
\makeatletter
\@ifundefined{lemma}{\newtheorem{lemma}{Lemma}}{}
\makeatother
\def \R{{\mathbb R}}
\title[Polyharmonic curves in semi-Riemannian manifolds]{Polyharmonic curves in semi-Riemannian manifolds}
\author{Stefano Montaldo, Andrea Ratto, Antonio Sanna}
\date{Source: arXiv:2506.16954v1 (CC BY 4.0)}
\begin{document}
\maketitle

\noindent\emph{Source and licence.} Polyharmonic curves in semi-Riemannian manifolds. Version arXiv:2506.16954v1 (CC BY 4.0), distributed under the Creative Commons Attribution 4.0 International licence, as recorded in the arXiv licence metadata for this exact version and shown on the abstract page https://arxiv.org/abs/2506.16954v1. Full rights notice: licenses/arxiv-2506.16954-CC-BY-4.0.txt.\par\smallskip

\noindent\textit{Editorial scope.} Counted unit: the Lemma whose source label is \texttt{lemma-gamma-elica} in the article (source0.tex lines 967--980, source label \texttt{lemma-gamma-elica}), delivered together with the article's own complete original proof of it (source0.tex lines 981--1017, one \texttt{proof} environment of 37 lines). No mathematics is written, repaired or re-derived: statement and proof are the article's own text line for line. Required context retained uncounted: eq-r-harm-general-helices-r>=4 (lines 407--409); eq-d-condizione (lines 427--429). Every definition, display and label of those lines is retained unchanged. Omitted from this package and not redistributed: 1--406, 410--426, 430--966, 1018--1436. Nothing inside a retained excerpt is omitted or altered: every retained body is delivered exactly as the article's lines are written, so no omission receipt of any kind accompanies this package. Citations: the retained excerpts carry no citation command at all, so no bibliography entry of the article is transcribed and no reference is redistributed. Reference adaptation: none was needed and none was made; every reference inside the delivered unit resolves inside it, so no unresolved reference or citation is produced. Printed slips kept literal: no printed word, symbol, hypothesis or number of the counted statement or of its proof was altered. Layout and class: the article's own class is \texttt{amsart} with packages including amsfonts, amscd, amssymb, url, textcomp, babel, color, hyperref; the wrapper is the lane's pinned \texttt{amsart} editorial wrapper at 10pt on A4, which restates the theorem-like environments the retained text needs and adds the article's own macro definitions that the retained text needs (\texttt{\textbackslash R}), with extra wrapper packages (bm, bbm, dsfont, xspace, cancel, mathtools). The delivered numbering is the unit's own. Assets: no retained span contains a figure environment, a graphics-inclusion command, a PDF-page inclusion, a table or a TikZ picture; no image file is copied into this package, the package declares no asset dependency and no image file is redistributed with it. Licence: version arXiv:2506.16954v1 of this article is distributed under the Creative Commons Attribution 4.0 International licence, as recorded in the arXiv licence metadata for this exact version and shown on the abstract page https://arxiv.org/abs/2506.16954v1. A full rights notice is delivered at licenses/arxiv-2506.16954-CC-BY-4.0.txt. Characters: every retained excerpt is pure ASCII, so the delivered engine, XeLaTeX with its default Latin Modern text font, reports no missing character.
\begin{equation}\label{eq-r-harm-general-helices-r>=4}
\epsilon_1 \kappa^2+\epsilon_3 \tau^2=0\,,
\end{equation}

\begin{equation}\label{eq-d-condizione}
\epsilon_3 \left (\tau_\alpha^2-\epsilon_1 \,d_1^2\,\kappa_\alpha^2 \right )>0\,,
\end{equation}

\begin{lemma}\label{lemma-gamma-elica} The curve $\gamma$ verifies the Frenet equations
\begin{equation}\label{Fren-eq-pseudo-bis}
\begin{cases}
\nabla_T T=\epsilon_2 \kappa N \\
\nabla_T N=-\epsilon_1 \kappa T+\epsilon_3 \tau B\\
\nabla_T B=-\epsilon_2 \tau N\,,
\end{cases}
\end{equation}
with $\epsilon_2=1$ and
\begin{equation}\label{eq-kappaetauproduct}
\kappa=(\epsilon_1+d_1^2)\kappa_\alpha\,, \quad \quad \tau=\sqrt{\epsilon_3\,(\epsilon_1+d_1^2) \left (\tau_\alpha^2-\epsilon_1 \,d_1^2\,\kappa_\alpha^2 \right )}\,.
\end{equation}
Moreover, $\gamma$ does not verify \eqref{eq-r-harm-general-helices-r>=4}.
\end{lemma}

\begin{proof} We compute
\[
\nabla_T T= \nabla_{ d_1\, \partial_t+d_2\, T_\alpha} \left ( d_1\, \partial_t+d_2\, T_\alpha\right )=d_2^2 \widetilde{\nabla}_{T_\alpha} T_{\alpha}= d_2^2 \kappa_\alpha N_\alpha
\]
from which we deduce $N=N_\alpha$, $\epsilon_2=1$ and $\kappa$ as in \eqref{eq-kappaetauproduct}. Similarly
\[
\nabla_T N=d_2 \widetilde{\nabla}_{T_\alpha} N_\alpha=-d_2\, \kappa_\alpha T_\alpha + d_2\, \tau_\alpha B_\alpha\,.
\]
Next, equalling the right-hand side to $-\epsilon_1 \kappa T + \epsilon_3 \tau B$ and considering the components tangent to $\R$ and $N^{m-1}(c)$ respectively we find:
\begin{equation*}\label{eq-1-lemma}
\begin{array}{ll}
{\rm (i)}& \epsilon_1 \kappa d_1 + \epsilon_3 \tau B_t=0 \\
{\rm (ii)}& \big (\epsilon_1 \kappa d_2-\kappa/d_2 \big )T_\alpha + d_2 \tau_\alpha B_\alpha-\epsilon_3 \tau \widetilde{B}=0
\end{array}
\end{equation*}
which imply
\begin{equation}\label{eq-2-lemma}
\begin{array}{ll}
{\rm (i)}& B_t=-\epsilon_1 \epsilon_3 (\kappa / \tau) d_1 \\
{\rm (ii)}& \widetilde{B}=(\epsilon_3 /  \tau)\Big (\big (\epsilon_1 \kappa d_2-\kappa/d_2 \big )T_\alpha + d_2\tau_\alpha B_\alpha\Big )
\end{array}
\end{equation}
We use these facts to compute the expression of $\tau$:
\begin{eqnarray}\nonumber
\tau&=& \langle \nabla_T N,B \rangle \\\nonumber
&=&\langle -d_2\, \kappa_\alpha T_\alpha + d_2\, \tau_\alpha B_\alpha,-B_t \,\partial_t + \widetilde{B} \rangle\\
&=&\langle -d_2\, \kappa_\alpha T_\alpha + d_2\, \tau_\alpha B_\alpha,\frac{\epsilon_3}{\tau} \Big (\big (\epsilon_1 \kappa d_2-\kappa/d_2 \big )T_\alpha + d_2 \tau_\alpha B_\alpha\Big )\rangle\\\nonumber
&=&\frac{\epsilon_3}{\tau} \Big (d_2^2 \tau^2_\alpha -\kappa^2(\epsilon_1-\frac{1}{d_2^2} \Big )\,.
\end{eqnarray}
Next, using $d_2^2=\epsilon_1+d_1^2$ and the assumption \eqref{eq-d-condizione} we easily deduce the expression for $\tau$ given in \eqref{eq-kappaetauproduct}.

Finally, a simple computation yields
\[
\epsilon_1 \kappa^2 +\epsilon_3 \tau^2 =\left(d_1^2+\epsilon_1 \right)\left(\kappa^2_\alpha +\tau^2_\alpha \right )\neq 0\,.
\]
Thus $\gamma$ does not verify condition \eqref{eq-r-harm-general-helices-r>=4}, as required to end the lemma.
\end{proof}

\end{document}
